% Independently accepted and integrated into sections/C.tex on 18 November 2026.
\begin{claim}{C.25}{Proved under stated hypotheses}{A structural obstruction to a tensor on the square-zero example}\label{C25}
Let $\mathcal A$ be an abelian category with an additive monoidal product and unit $1$. Suppose
\[
0\longrightarrow S\xrightarrow{i}D\xrightarrow{p}1\longrightarrow0
\]
is exact. For every $X\in\mathcal A$, the map
\[
\operatorname{Hom}_{\mathcal A}(1,S\otimes X)
\longrightarrow\operatorname{Hom}_{\mathcal A}(1,D\otimes X)
\]
induced by $i\otimes X$ is injective. We require neither exactness of the tensor nor a symmetry.

Consequently the category in Claim~\ref{C24} admits no $k$-bilinear monoidal product with unit $\kappa k$ for which the canonical scalar-copower comparisons
\[
W\otimes_k S\longrightarrow S\otimes\kappa W,
\qquad
W\otimes_k D\longrightarrow D\otimes\kappa W
\]
are isomorphisms for every vector space $W$, where $S=P_1(k)$ and $D=P_0(k)$. In particular it admits no $k$-bilinear monoidal product with that unit which preserves arbitrary coproducts in its second variable. This excludes every cocontinuous symmetric $k$-bilinear tensor on that example. It does not exclude a merely bilinear tensor without the indicated compatibility.
\end{claim}
\begin{proof}
We suppress the associativity and unit isomorphisms; the following equations are invariant under their insertion by monoidal coherence. Define
\[
r_1=p\otimes\id_D:D\otimes D\longrightarrow D,
\qquad
r_2=\id_D\otimes p:D\otimes D\longrightarrow D.
\]
The interchange and unit laws give $pr_1=p\otimes p=pr_2$. Since $i$ is the kernel of $p$, there is a unique
\[
d:D\otimes D\longrightarrow S,
\qquad i\,d=r_1-r_2.
\]
After precomposition with $\id_D\otimes i$, we obtain
\[
i\,d(\id_D\otimes i)
=i(p\otimes\id_S)-(\id_D\otimes pi)
=i(p\otimes\id_S).
\]
Cancellation of the monomorphism $i$ therefore gives
\[
d(\id_D\otimes i)=p\otimes\id_S.
\]
Let $\beta:1\to S\otimes X$ satisfy $(i\otimes\id_X)\beta=0$. Tensoring this equality on the left by $D$ and then composing with $d\otimes\id_X$, we find
\[
0=(d\otimes\id_X)(\id_D\otimes i\otimes\id_X)
  (\id_D\otimes\beta)
 =(p\otimes\id_{S\otimes X})(\id_D\otimes\beta)
 =\beta p.
\]
Because $p$ is an epimorphism, $\beta=0$. This proves the general assertion without supposing that $i\otimes X$ itself is monic.

Now use the category and notation of Claim~\ref{C24}. Since $E(k)=0$, its first short exact sequence at $k$ is
\[
0\longrightarrow P_1(k)\xrightarrow{i}P_0(k)
\xrightarrow{p}\kappa k\longrightarrow0.
\]
The actual scalar copowers of $S=P_1(k)$ and $D=P_0(k)$ are $P_1(W)$ and $P_0(W)$: this follows either from the free adjunction or from their displayed Hom formulas. Under these identifications the copower of $i$ is
\[
i_W:P_1(W)\longrightarrow P_0(W),\qquad i_W=(0,\id_W),
\]
and its kernel is $\kappa E(W)$. Choose an infinite-dimensional $W$ and a nonzero $e\in E(W)$. The corresponding map
\[
\beta_e:\kappa k\longrightarrow\kappa E(W)
\longrightarrow P_1(W)
\]
is nonzero and is killed by $i_W$. The canonical copower comparisons are natural in their object argument. If they are isomorphisms for both $S$ and $D$, they therefore identify $i_W$ with $i\otimes\kappa W$. The first part, applied with $X=\kappa W$, contradicts the existence of $\beta_e$.

Finally, if the second-variable tensor functors preserve coproducts, write $W$ as a free vector space on a basis. Since $\kappa W$ is the corresponding actual coproduct of units, the two canonical comparisons are isomorphisms. The contradiction applies. No choice of basis enters the naturality of those canonical maps, and no preservation of cokernels is needed for this last implication over a field.
\end{proof}
